% generated from Paper 10/anccft10.tex
\chapter[\texorpdfstring{X. $\Atwo$ complexes, higher simplicial Laplacians, and the fate of the trace obstruction in
rank two}{X. A2 complexes, higher simplicial Laplacians, and the fate of the trace obstruction in rank two}]{\texorpdfstring{Algorithmic non-commutative class field theory, X: $\Atwo$ complexes, higher simplicial Laplacians, and the fate of the trace obstruction in
rank two}{Algorithmic non-commutative class field theory, X: A2 complexes, higher simplicial Laplacians, and the fate of the trace obstruction in rank two}}\label{chap:X}
\chaptermark{X}
\begin{chapterabstract}
The series has so far lived on quotients of the Bruhat--Tits tree of $\PGL_2$. This paper
moves one rank up, to finite $\Atwo$ complexes --- quotients of two-dimensional buildings of
order $q$, locally the building of $\PGL_3$ over a local field with residue field $\F_q$ ---
and asks which of the rank-one theorems survive. Three things are established. First, the
Hodge Laplacians $\Delta_0,\Delta_1,\Delta_2$ of the typed chain complex give higher
critical groups $\Jac_i(Y)=\Tor\coker\Delta_i$, computed exactly on thin tori and on five
thick complexes of order $2$ with $21$, $24$, $42$, $84$ and $168$ vertices, built here from
a triangle presentation over $\PG(2,2)$ compatible with a Singer cycle and verified to have
Heawood-graph links; the Hodge vertex Laplacian is no longer a Hecke operator, and the
arithmetic bridge of [Ch.~\ref{chap:I}, Thm.~3.5] is carried instead by the \emph{Hecke--Laplacian}
$\DeltaH=(q^{3}-1)I+A_1-qA_2$, the value at $u=1$ of the cubic pencil
$P(u)=\det(I-uA_1+qu^{2}A_2-q^{3}u^{3}I)$, whose critical group has order
$\frac1n\prod_j|P_j(1)|$, a product of inverse local $\mathrm{GL}_3$ $L$-values. Second,
the companion $\Cthree$ of the cubic pencil is unimodularly equivalent to $\DeltaH\oplus I\oplus I$
exactly as in [Ch.~\ref{chap:II}], but the trace obstruction of [Ch.~\ref{chap:II}] does \emph{not} transfer: the $\Z/3$
type grading forces $\Tr(\Cthree)^{k}=0$ for $3\nmid k$, and a block computation gives
$\Tr(\Cthree)^{3}=n(q-1)^{2}(q+1)\ge0$. There is no cubic trace obstruction; on the five
thick complexes, all of which turn out to be Ramanujan, every net trace of the reduced
(type-preserving) pencil is non-negative and its Perron condition holds, so by the theorem
of Kim--Ormes--Roush the nonzero spectrum of $\Cthree$ is that of a non-negative integer
matrix, irreducible of period three --- the rank-two pencil is spectrally realizable, the
opposite of the rank-one situation, and the realization problem for the $\PGL_3$ Hecke
critical group is open and unobstructed. Third, the shadow construction of [Ch.~\ref{chap:VI}] extends to
arbitrary
integer matrices with positive diagonal by a signed-shadow lemma in which positive
off-diagonal entries are produced by one-directional auxiliary vertices with double loops;
it yields unital Kirchberg algebras $\cO(Y,\Delta_1)$ and $\cO(Y,\DeltaH)$ with
$K_0\cong\Z^{b_1}\oplus\Jac_1(Y)$ and $K_0\cong\Z\oplus\JacH(Y)$ respectively, functorially
in the complex. What the specification predicted and what is proved differ in two places,
and the paper says where.
\end{chapterabstract}

\section*{Status of the results}

Every numbered statement is proved. The inputs are finite linear algebra, the classical
Cuntz--Krieger facts quoted in [Chs.~\ref{chap:VI}, \ref{chap:VIII}], the Cartwright--Mantero--Steger--Zappa theorem
that a triangle presentation gives a group acting simply transitively on the vertices of an
$\Atwo$ building \cite{CMSZ}, and, for one corollary that nothing else depends on, the
Kim--Ormes--Roush theorem on the nonzero spectra of primitive integer matrices \cite{KOR}.
The Kang--Li zeta function of $\PGL_3$ complexes \cite{KL} is context for the cubic pencil,
not an input. Every finite identity was machine-verified exactly (Table~\ref{X:tab:battery});
the five thick complexes are constructed by the script, and their local structure (Heawood
links, regularity, flag condition) is checked rather than assumed --- with the result that
the flag condition holds only for the largest of them, a distinction the theorems respect. Section~\ref{X:sec:scope}
records what is open and the corrections to the specification.

\section{\texorpdfstring{$\Atwo$ complexes of order $q$}{A2 complexes of order q}}\label{X:sec:complexes}

\begin{definition}[$\Atwo$ complex]\label{X:def:a2}
An \emph{$\Atwo$ complex of order $q\ge1$} is a finite connected pure two-dimensional
simplicial complex $Y$ with vertex set $V$ ($n:=|V|$), edge set $E$ and triangle set $F$,
together with a type function $\tau\colon V\to\Z/3$, such that
\begin{enumerate}
\item[(i)] every triangle has one vertex of each type;
\item[(ii)] every vertex has exactly $q^{2}+q+1$ neighbours of type $\tau(v)+1$ and
$q^{2}+q+1$ of type $\tau(v)+2$;
\item[(iii)] every edge lies in exactly $q+1$ triangles;
\item[(iv)] (flag condition) every closed walk $x\to y\to z\to x$ with
$\tau(y)=\tau(x)+1$, $\tau(z)=\tau(x)+2$ is the boundary of a triangle.
\end{enumerate}
Every edge is oriented from type $t$ to type $t+1$ and every triangle $(a,b,c)$ is listed
with types $(t,t+1,t+2)$, so that $\partial_2(a,b,c)=(a,b)+(b,c)+(c,a)$ and
$\partial_1(a,b)=b-a$ are canonical; $A_1\in M_n(\{0,1\})$ is the adjacency $x\to y$ for
edges, $A_2:=A_1^{t}$, and we assume $A_1A_2=A_2A_1$ (true for building quotients, where
$A_1,A_2$ generate the commutative spherical Hecke algebra; verified on every complex used).
\end{definition}

Quotients of a thick $\Atwo$ building of order $q\ge2$ by a type-rotating, freely acting
group with injectivity radius $\ge2$ are $\Atwo$ complexes in this sense, with vertex links
the incidence graph of $\PG(2,q)$; the flag condition holds because buildings are flag
complexes. Thin complexes ($q=1$) are triangulated flat tori (Construction~\ref{X:con:thin}).

\begin{lemma}[Counts]\label{X:lem:counts}
$|E|=n(q^{2}+q+1)$, $|F|=\tfrac13n(q^{2}+q+1)(q+1)$, $\Tr A_1^{3}=3|F|$,
$\Tr(A_1A_2)=|E|$, and $\chi(Y)=n\bigl(1+\tfrac13(q^{2}+q+1)(q-2)\bigr)$. For $q=2$:
$|E|=|F|=7n$ and $\chi(Y)=n$.
\end{lemma}

\begin{proof}
(ii) gives $|E|$; each vertex lies in $(q^{2}+q+1)(q+1)$ triangles by (ii),(iii), each
counted three times; $\Tr A_1^{3}$ counts closed type-increasing $3$-walks, three per
triangle by (iv); $\Tr(A_1A_1^{t})=\sum_{x,y}A_1(x,y)^{2}=|E|$.
\end{proof}

\begin{construction}[Thin tori]\label{X:con:thin}
On $\Z^{2}$ let $d_1=(1,0)$, $d_2=(0,1)$, $d_3=(-1,-1)$, edges $v\to v+d_i$, type
$\tau(a,b)=a+b\bmod3$, triangles $\{v,v+d_i,v+d_i+d_j\}$, $i\neq j$. This is the
$\Atwo$ Coxeter complex (one apartment); its quotient by $m\Z^{2}$, $3\mid m$, $m\ge3$, is an
$\Atwo$ complex of order $1$ with $m^{2}$ vertices, $b_1=2$, $b_2=1$, hexagonal links.
\end{construction}

\begin{construction}[Thick complexes of order $2$]\label{X:con:thick}
Label the points of $\PG(2,2)$ by $\Z/7$ with lines $\{i,i+1,i+3\}$ and let
$\lambda(x)=\{x+1,x+2,x+4\}$, the Singer cycle. A triangle presentation compatible with
$\lambda$ \cite{CMSZ} is a set $T$ of $21$ triples with (A) $(x,y,z)\in T\Rightarrow
(y,z,x)\in T$, (B) $(x,y,\cdot)\in T$ iff $y\in\lambda(x)$, (C) the third entry is
determined by the first two. Exhaustive search over all bijections without absolute points
finds, for $\lambda$ the Singer cycle, the presentation with cyclic orbits
\[
T/\Z_3=\{(0,1,3),(0,2,6),(0,4,5),(1,2,4),(1,5,6),(2,3,5),(3,4,6)\}
\]
(axioms re-verified by the script). By \cite{CMSZ} the group
$\Gamma=\langle a_0,\dots,a_6\mid a_xa_ya_z=1,\ (x,y,z)\in T\rangle$ acts simply
transitively and type-rotatingly on the vertices of an $\Atwo$ building of order $2$; its
type-preserving subgroup $\Gamma^{0}$ has index $3$, and $X_3:=\Gamma^{0}\backslash B$ is
the $\Delta$-complex with three vertices, $21$ edges and $21$ triangles. We compute
$H_1(X_3;\Z)\cong\Z/2\oplus\Z/2\oplus\Z/14$ and build the abelian voltage covers of $X_3$
with deck group $A$ for the quotients $A=\Z/7$, $(\Z/2)^{3}$, $\Z/14$ and $H_1(X_3)$
itself, exactly as the derived graphs of [Ch.~\ref{chap:II}, \S5] in one dimension higher: vertices
$(i,g)$, edges lifted by voltages in $A$, triangles lifting because their boundaries are
null-homologous. The results $Y_{21},Y_{24},Y_{42},Y_{84},Y_{168}$ are simplicial, every
vertex link is the Heawood graph (bipartite, $3$-regular, $14$ vertices, girth $6$), and
conditions (i)--(iii) hold, with $|E|=|F|=7n$ (Table~\ref{X:tab:battery}, line (P)). The flag
condition (iv) holds for $Y_{168}$ and fails for the four smaller covers: there the deck
group is too small, some product $a_xa_ya_{z'}$ with $(x,y,z')\notin T$ has trivial voltage,
and $\Tr A_1^{3}$ exceeds $3|F|$ by $28n$, $21n$, $12n$, $4n$ respectively. So (iv) is not
implied by the link condition, and the five complexes separate the two hypotheses cleanly:
$Y_{168}$ is an $\Atwo$ complex in the full sense of Definition~\ref{X:def:a2}, the others
satisfy (i)--(iii) only. The count $168=|\PGL_3(\F_2)|$ for the maximal abelian cover
suggests that $Y_{168}$ is a congruence quotient; we neither need nor claim this. Whether the
building is that of $\PGL_3(\F_2((t)))$, of $\PGL_3(\Q_2)$, or exotic is not decided here:
every statement below depends only on the local structure of Definition~\ref{X:def:a2}, and
the statements that use (iv) are marked.
\end{construction}

\section{Hodge Laplacians and higher critical groups}\label{X:sec:hodge}

\begin{definition}\label{X:def:hodge}
$\Delta_0:=\partial_1\partial_1^{t}$ on $C_0$, $\Delta_1:=\partial_1^{t}\partial_1
+\partial_2\partial_2^{t}$ on $C_1$, $\Delta_2:=\partial_2^{t}\partial_2$ on $C_2$, and
$\Jac_i(Y):=\Tor\coker_\Z\Delta_i$, the \emph{$i$-th critical group}. Following
Duval--Klivans--Martin \cite{DKM} we also record $K(Y):=\ker\partial_1/\operatorname{im}
(\partial_2\partial_2^{t})$.
\end{definition}

\begin{proposition}[Structure]\label{X:prop:hodge}
$\partial_1\partial_2=0$; $\coker\Delta_i\cong\Z^{b_i}\oplus\Jac_i(Y)$ with $b_i$ the Betti
numbers of $Y$; $b_0=1$; $\Delta_0=2(q^{2}+q+1)I-(A_1+A_2)$ is the Laplacian of the
$1$-skeleton, so $\Jac_0(Y)$ is its sandpile group with $|\Jac_0(Y)|$ the number of
spanning trees of the $1$-skeleton; $\Delta_1$ has diagonal entries $q+3$, entry $0$ at pairs
of edges spanning a triangle, and entries $\pm1$ at pairs sharing a vertex and no triangle.
\end{proposition}

\begin{proof}
$\partial_1\partial_2(a,b,c)=(b-a)+(c-b)+(a-c)=0$. Each $\Delta_i$ is symmetric with
$\ker\Delta_i=\ker\partial_i\cap\ker\partial_{i+1}^{t}$ of rank $b_i$ (Hodge), so its
cokernel has free rank $b_i$. The diagonal of $\Delta_1$ is $2+(q+1)$ by (iii); for edges
$e=(a,b)$, $f=(b,c)$ in a triangle the down-term $\langle\partial e,\partial f\rangle=-1$
cancels the up-term $+1$, and without a triangle only the down-term survives.
\end{proof}

\section{The Hecke--Laplacian and the \texorpdfstring{$\mathrm{GL}_3$ $L$-values}{GL3 L-values}}\label{X:sec:hecke}

\begin{definition}[Cubic pencil and Hecke--Laplacian]\label{X:def:pencil}
$P(u):=I-uA_1+qu^{2}A_2-q^{3}u^{3}I\in M_n(\Z[u])$ and
\[
\DeltaH:=-P(1)=(q^{3}-1)I+A_1-qA_2,\qquad \JacH(Y):=\Tor\coker_\Z\DeltaH .
\]
For a joint eigenpair $(a_j,b_j)$ of $(A_1,A_2)$ write $P_j(u)=1-a_ju+qb_ju^{2}-q^{3}u^{3}
=\prod_{i=1}^{3}(1-\lambda_{j,i}u)$; the trivial pair $a=b=q^{2}+q+1$ gives
$(1-u)(1-qu)(1-q^{2}u)$.
\end{definition}

The pencil is the vertex factor of the Kang--Li zeta function of a $\PGL_3$ complex
\cite{KL}, and $P_j(u)^{-1}$ is the local $L$-factor of the unramified $\mathrm{GL}_3$
Hecke eigenform with Satake parameters $\lambda_{j,i}/q$; we use only the polynomial.

\begin{theorem}[Hecke critical group and $L$-values]\label{X:thm:hecke}
\begin{enumerate}
\item[(i)] $\DeltaH\mathbf1=0$ and $\mathbf1^{t}\DeltaH=0$.
\item[(ii)] If $P_j(1)\neq0$ for every nontrivial $j$ (no nontrivial eigenform has a
Satake parameter $\lambda_{j,i}=1$; automatic when $|\lambda_{j,i}|=q$), then
$\rk\DeltaH=n-1$, all $n$ principal cofactors of $\DeltaH$ are equal, and
\[
\coker_\Z\DeltaH\cong\Z\oplus\JacH(Y),\qquad
|\JacH(Y)|=\frac1n\prod_{j\ \mathrm{nontrivial}}\bigl|P_j(1)\bigr| .
\]
\item[(iii)] $\DeltaH\neq\Delta_0$ for $q\ge1$, and $\Delta_0$ is not a specialization of the
pencil: no $u$ gives $P(u)\in\Q^{\times}\Delta_0$. For $\PGL_2$ the two operators coincide
($(q{+}1)I-T_p=\Delta_0=-P(1)$ with $P(u)=I-uT_p+qu^{2}$).
\end{enumerate}
\end{theorem}

\begin{proof}
(i) Row sums: $(q^{3}-1)+(q^{2}+q+1)-q(q^{2}+q+1)=0$; column sums likewise since
$\mathbf1^{t}A_1=\mathbf1^{t}A_2=(q^{2}+q+1)\mathbf1^{t}$. (ii) The eigenvalues of $\DeltaH$
on the joint eigenspaces are $-P_j(1)$, zero exactly for the trivial pair under the
hypothesis, so the kernel is $\Z\mathbf1$ on both sides; the equal-cofactor argument of
[Ch.~\ref{chap:III}, proof of Thm.~2.1, Step~2] applies verbatim to a matrix with two-sided kernel
$\mathbf1$ and gives $\prod_{\mu\neq0}\mu=n\kappa$ with $\kappa$ the common cofactor
$=|\Tor\coker|$. (iii) $P(u)\in\Q\Delta_0$ would need $-u:qu^{2}=-1:-1$, i.e.\ $u=-1/q$,
where $P(-1/q)=2I+(A_1+A_2)/q$ has the wrong constant term unless $2q=2(q^{2}+q+1)$,
impossible.
\end{proof}

\begin{theorem}[Companion reduction]\label{X:thm:companion}
Let $\Cthree=\begin{psmallmatrix}A_1&-qA_2&q^{3}I\\ I&0&0\\0&I&0\end{psmallmatrix}$, so
that $\det(I-u\Cthree)=\det P(u)$. With
$U=\begin{psmallmatrix}I&q^{3}I-qA_2&q^{3}I\\0&I&0\\0&0&I\end{psmallmatrix}$ and
$V=\begin{psmallmatrix}I&0&0\\ I&I&0\\ I&I&I\end{psmallmatrix}$, both unimodular,
\[
U\,(I-\Cthree)\,V=P(1)\oplus I_n\oplus I_n=(-\DeltaH)\oplus I_n\oplus I_n,
\qquad\text{hence}\qquad \coker_\Z(I-\Cthree)\cong\Z\oplus\JacH(Y).
\]
\end{theorem}

\begin{proof}
$I-\Cthree=\begin{psmallmatrix}I-A_1&qA_2&-q^{3}I\\-I&I&0\\0&-I&I\end{psmallmatrix}$; the
first block row of $U(I-\Cthree)$ is $(I-A_1,qA_2,-q^{3}I)+(q^{3}I-qA_2)(-I,I,0)
+q^{3}(0,-I,I)=(P(1),0,0)$, and the column operations $V$ clear the subdiagonal $-I$'s.
This is [Ch.~\ref{chap:II}, \S1.1] one degree higher.
\end{proof}

\section{The fate of the trace obstruction}\label{X:sec:trace}

\begin{theorem}[Grading]\label{X:thm:grading}
$\Tr(A_1^{\alpha}A_2^{\beta})=0$ unless $\alpha\equiv\beta\pmod3$. Consequently
$\Tr(\Cthree)^{k}=0$ for every $k$ with $3\nmid k$, and the spectrum of $\Cthree$ is
invariant under multiplication by cube roots of unity.
\end{theorem}

\begin{proof}
A closed walk contributing to the trace changes type by $\alpha+2\beta\equiv\alpha-\beta$.
$\Tr(\Cthree)^{k}=\sum_jp_k(a_j,b_j)$ with $p_k$ the $k$-th power sum of the roots of
$\lambda^{3}-a\lambda^{2}+qb\lambda-q^{3}$, a polynomial in monomials $a^{\alpha}(qb)^{\beta}
(q^{3})^{\gamma}$ with $\alpha+2\beta+3\gamma=k$, so $\alpha\equiv\beta$ forces $3\mid k$.
Equivalently $D_\omega=\operatorname{diag}(\omega^{\tau(v)})\oplus\cdots$ conjugates
$\omega\Cthree$ to $\Cthree$.
\end{proof}

\begin{theorem}[No cubic trace obstruction]\label{X:thm:notrace}
For every $\Atwo$ complex of order $q$,
\[
\Tr\Cthree=0,\qquad\Tr(\Cthree)^{2}=0,\qquad
\Tr(\Cthree)^{3}=3|F|-3q|E|+3nq^{3}=n(q-1)^{2}(q+1)\ \ge\ 0,
\]
with equality only in the thin case $q=1$. Hence the mechanism of {\rm[Ch.~\ref{chap:II}, Thms.~1.5--1.6]}
--- a negative trace of a power of the companion, incompatible with any non-negative
matrix of any size --- has no analogue in rank two at orders $\le5$: the first net trace
that is not forced to be non-negative by these identities is $\tr_6$.
\end{theorem}

\begin{proof}
$p_1=a$, $p_2=a^{2}-2qb$, $p_3=a^{3}-3qab+3q^{3}$ (Newton). $\Tr A_1=0$ (no loops),
$\Tr A_1^{2}=0$ and $\Tr A_2=0$ by Theorem~\ref{X:thm:grading}, $\Tr A_1^{3}=3|F|$ and
$\Tr(A_1A_2)=|E|$ by Lemma~\ref{X:lem:counts}; substitute the counts:
$(q^{2}+q+1)(q+1)-3q(q^{2}+q+1)+3q^{3}=q^{3}-q^{2}-q+1=(q-1)^{2}(q+1)$. The net trace
$\tr_m=\sum_{d\mid m}\mu(m/d)\Tr(\Cthree)^{d}$ vanishes for $m\le5$, $m\neq3$, and equals
$\Tr(\Cthree)^{3}\ge0$ for $m=3$. Without the flag condition the identity reads
$\Tr(\Cthree)^{3}=\Tr A_1^{3}-3q|E|+3nq^{3}$, still non-negative since $\Tr A_1^{3}\ge3|F|$:
on $Y_{168}$ the value is $504=n(q-1)^{2}(q+1)$ exactly; on $Y_{21}$ it is $651=63+588$,
the excess being the spurious closed $3$-walks (Table~\ref{X:tab:battery}).
\end{proof}

\begin{remark}[Why rank one is different]\label{X:rem:why}
In [Ch.~\ref{chap:II}] the companion of $1-uT_p+qu^{2}$ has $\Tr C^{2}=\Tr T_p^{2}-2qn=n(q+1)-2qn<0$:
the $\Z/2$ type structure of a graph allows a walk to return in two steps (edge reversal),
so the symmetric adjacency contributes $\Tr T_p^{2}=2|E|>0$ but not enough to beat the
$-2qn$ of the negative block. In rank two the type grading is $\Z/3$: a type-increasing walk
cannot return before three steps, the negative block $-qA_2$ is invisible to
$\Tr(\Cthree)^{k}$ for $k<3$, and at $k=3$ the chambers contribute $3|F|=n(q^{2}+q+1)(q+1)$,
which exceeds the $3q|E|-3nq^{3}=n(q^{2}+q+1)\cdot3q-3nq^{3}$ removed by the negative block.
The obstruction of [Ch.~\ref{chap:II}] is a rank-one accident of parity.
\end{remark}

\begin{definition}[Reduced pencil]\label{X:def:reduced}
By Theorem~\ref{X:thm:grading}, $\det P(u)\in\Z[u^{3}]$; write $\det P(u)=Q(u^{3})$ with
$Q\in\Z[v]$ of degree $n$. Its roots are $\lambda^{3}$, one for each $\omega$-orbit
$\{\lambda,\omega\lambda,\omega^{2}\lambda\}$ of eigenvalues of $\Cthree$: the trivial orbits
give $1,q^{3},q^{6}$, and the power sums of the multiset $\Lambda_Q$ of roots are
$p_d(\Lambda_Q)=\tfrac13\Tr(\Cthree)^{3d}$. $\Lambda_Q$ is the spectrum of the
type-preserving (three-step) dynamics.
\end{definition}

\begin{theorem}[Spectral realizability of the thick examples]\label{X:thm:realizable}
For each of $Y_{21},Y_{24},Y_{42},Y_{84},Y_{168}$:
\begin{enumerate}
\item[(i)] $\Cthree$ has exactly nine eigenvalues of modulus $\ge q$: the orbits of $1,q,q^{2}$
under $\omega$; every other eigenvalue has modulus exactly $q$ (numeric, to $10^{-6}$): the
complexes are Ramanujan in the sense of \cite{KL}. In particular $\Cthree$ itself never
satisfies the Perron condition of a primitive matrix --- $q^{2},q^{2}\omega,q^{2}\omega^{2}$
are all dominant --- which is forced by the grading and is the signature of an irreducible
matrix of period $3$.
\item[(ii)] For $\Lambda_Q$: $q^{6}$ is simple and strictly dominant; the net traces
\[
\tr_m(\Lambda_Q)=\sum_{d\mid m}\mu(m/d)\,\tfrac13\Tr(\Cthree)^{3d}
\]
are non-negative for $m\le10$ (exact integers) and provably positive for
$m\ge m_0\in\{2,3\}$ (Lemma~\ref{X:lem:bound}); hence by \cite{KOR} there is a primitive
non-negative integer matrix $B_0$ with nonzero spectrum $\Lambda_Q$.
\item[(iii)] Consequently the block-cyclic non-negative integer matrix
$\begin{psmallmatrix}0&B_0&0\\0&0&I\\I&0&0\end{psmallmatrix}$ has the nonzero spectrum of
$\Cthree$: the $\PGL_3$ pencil of these complexes is spectrally realizable by a non-negative
integer matrix, irreducible of period $3$.
\end{enumerate}
\end{theorem}

\begin{proof}
(i) is computation (Table~\ref{X:tab:battery}, line (R)); the grading statement is
Theorem~\ref{X:thm:grading}. (ii) The Boyle--Handelman conditions for a primitive integer
matrix are: $\Lambda_Q$ is the root multiset of a monic integer polynomial ($Q$), a simple
strictly dominant positive root ($q^{6}$, by (i)), and $\tr_m(\Lambda_Q)\ge0$ for all $m$;
\cite{KOR} proves them sufficient. (iii) If $M=\begin{psmallmatrix}0&X&0\\0&0&Y\\Z&0&0
\end{psmallmatrix}$ then $M^{3}=\operatorname{diag}(XYZ,YZX,ZXY)$, so the nonzero spectrum
of $M$ consists of all cube roots of the nonzero spectrum of $XYZ=B_0$, i.e.\ of the
$\omega$-orbits of $\Lambda_Q^{1/3}$, which is the nonzero spectrum of $\Cthree$.
\end{proof}

\begin{lemma}[Tail bound]\label{X:lem:bound}
Let $M<q^{2}$ bound the moduli of the genuinely nontrivial eigenvalues and $N=n-1$. Then
$\tr_m(\Lambda_Q)\ge q^{6m}+q^{3m}+1-NM^{3m}-\sum_{d\le m/2}\bigl(q^{6d}+q^{3d}+1+NM^{3d}\bigr)$,
positive for all $m\ge m_0$ with $m_0$ computable from $(n,q,M)$; for the five complexes
$m_0\le3$.
\end{lemma}

\begin{proof}
$|p_d(\Lambda_Q)-(1+q^{3d}+q^{6d})|\le NM^{3d}$; the proper divisors of $m$ are $\le m/2$; the
bound increases with $m$ once $q^{6m}$ dominates, which for $M=q$ and $n\le168$ happens by
$m=3$.
\end{proof}

\begin{remark}[What remains of the realization problem]\label{X:rem:realization}
Theorem~\ref{X:thm:realizable} is about the nonzero spectrum. The realization problem proper
--- a non-negative integer matrix $B$ shift equivalent over $\Z$ to $\Cthree$, so that
$\coker(I-B)\cong\Z\oplus\JacH(Y)$ with the Hecke dynamics --- asks for more than spectral
data, and we do not solve it. What is established is that the one tool which closed it in
rank one, Theorem~[Ch.~\ref{chap:II}, 1.6], is unavailable in rank two, and that the necessary conditions
known to us are all satisfied. The problem is open and, as far as the trace method can see,
unobstructed. Whether the double-loop mechanism of Section~\ref{X:sec:shadow} can be upgraded
from a $K$-theoretic to a dynamical realization is the natural next question, and [Ch.~\ref{chap:II}, Thm.~1.6] no longer forbids it here.
\end{remark}

\section{Signed shadows and the higher Kirchberg algebras}\label{X:sec:shadow}

\begin{lemma}[Signed shadow lemma]\label{X:lem:signed}
Let $M\in M_N(\Z)$ with $M_{ee}\ge1$ for all $e$, and let $\mathcal P=\{(e,f):e\neq f,\
M_{ef}>0\}$. Define a non-negative integer matrix $B$ on the vertex set
$\{e\}\sqcup\{w_{ef}:(e,f)\in\mathcal P\}\sqcup\{e'\}$ by the arcs
\begin{itemize}
\item $|M_{ef}|$ arcs $f\to e$ for each $e\neq f$ with $M_{ef}<0$;
\item for $(e,f)\in\mathcal P$: $M_{ef}$ arcs $f\to w_{ef}$, one arc $w_{ef}\to e$, two loops at
$w_{ef}$;
\item for each $e$: one arc $e\to e'$, $M_{ee}-1$ arcs $e'\to e$, two loops at $e'$.
\end{itemize}
Then, with $X:=(I-B^{t})_{\text{main},\text{aux}}$, $Y:=(I-B^{t})_{\text{aux},\text{main}}$
and $U=\begin{psmallmatrix}I&X\\0&I\end{psmallmatrix}$,
$V=\begin{psmallmatrix}I&0\\Y&I\end{psmallmatrix}$,
\[
U\,(I-B^{t})\,V\;=\;M\oplus(-I_{|\mathcal P|+N}) .
\]
If the graph on $\{e\}$ with edges $\{e,f\}$ for $M_{ef}\neq0$ is connected, $B$ is
irreducible with double-loop vertices, and $\cO_B:=C^*(\text{digraph of }B)$ is a unital
Kirchberg algebra with $K_0(\cO_B)\cong\coker_\Z M$ and $K_1(\cO_B)\cong\ker_\Z M$.
\end{lemma}

\begin{proof}
The auxiliary block of $I-B^{t}$ is $-I$ (each auxiliary vertex carries exactly two loops and
no other arcs among auxiliaries), so the Schur complement of that block is $R+XY$ where
$R$ is the main block: diagonal $1$, entry $M_{ef}$ at negative positions, $0$ elsewhere.
The auxiliary $w_{ef}$ contributes $X_{e,w}Y_{w,f}=(-1)(-M_{ef})=M_{ef}$ at $(e,f)$ and
nothing else; the shadow $e'$ contributes $X_{e,e'}Y_{e',e}=(-(M_{ee}-1))(-1)=M_{ee}-1$ at
$(e,e)$. Hence $R+XY=M$, and the displayed identity is the block factorization
$\begin{psmallmatrix}R&X\\Y&-I\end{psmallmatrix}=U^{-1}\bigl((R+XY)\oplus(-I)\bigr)V^{-1}$.
Irreducibility: every arc $f\to e$ or $f\to w_{ef}\to e$ is matched by the reverse
connection through $M_{fe}$ or through the shared vertex structure of a connected support
graph, and every $e'$ communicates with $e$ in both directions; the double loops give the
witness of [Ch.~\ref{chap:VI}, Thm.~3.1]; $K$-theory is the Cuntz--Krieger formula transported by $U,V$.
For $M=\Delta_p$ of a graph ($M_{ef}\le0$ off the diagonal) $\mathcal P=\emptyset$ and the
construction is [Ch.~\ref{chap:VI}, Def.~1.1] with $c=M_{ee}-1$.
\end{proof}

\begin{theorem}[Higher shadow algebras]\label{X:thm:shadow}
For an $\Atwo$ complex $Y$ let $\cO(Y,\Delta_1):=\cO_{B(\Delta_1)}$ and
$\cO(Y,\DeltaH):=\cO_{B(\DeltaH)}$ be the algebras of Lemma~\ref{X:lem:signed} for $M=\Delta_1$
(diagonal $q+3\ge1$) and $M=\DeltaH$ (diagonal $q^{3}-1\ge1$ for $q\ge2$). Both are unital
Kirchberg algebras, hence traceless, and
\[
K_0\bigl(\cO(Y,\Delta_1)\bigr)\cong\Z^{b_1}\oplus\Jac_1(Y),\quad K_1\cong\Z^{b_1};\qquad
K_0\bigl(\cO(Y,\DeltaH)\bigr)\cong\Z\oplus\JacH(Y),\quad K_1\cong\Z ,
\]
the second under the hypothesis of Theorem~\ref{X:thm:hecke}(ii). Type- and
orientation-preserving automorphisms of $Y$ act compatibly. Changing the edge orientations
conjugates $\Delta_1$ by a diagonal $\pm1$ matrix and changes $B(\Delta_1)$, but not the
$K$-theory.
\end{theorem}

\begin{proof}
Lemma~\ref{X:lem:signed} with Proposition~\ref{X:prop:hodge} and Theorem~\ref{X:thm:hecke}; the
support graph of $\Delta_1$ is connected because two edges of a connected complex are joined
by a chain of edges sharing vertices, and pairs spanning a triangle (where the entry is
$0$) are joined through a third edge; the support of $\DeltaH$ is the $1$-skeleton.
\end{proof}

\begin{remark}[Digraph, not cellular algebra]\label{X:rem:digraph}
The specification asked for a ``two-dimensional shadow complex with higher self-loop faces''
and a ``higher graph or cellular $C^*$-algebra''. No such theory is needed: the invariant to
be realized is a cokernel, cokernels are unimodular-equivalence invariants, and
Lemma~\ref{X:lem:signed} realizes any integer matrix with positive diagonal as the
Bowen--Franks presentation of an ordinary graph algebra. The two dimensions of $Y$ enter
through $\Delta_1$ and $\DeltaH$, not through the algebra.
\end{remark}

\section{The verification battery}\label{X:sec:battery}

All computations are exact except the spectra. Smith forms are computed $p$-adically:
the determinant of the (reduced) matrix by Bareiss elimination, its prime factors by trial
division, and the $p$-primary invariants by elimination over $\Z/p^{k}$ with
valuation-minimal pivots; for singular symmetric $M$ the torsion found is certified
complete by the lattice identity
\begin{equation}\label{X:eq:lattice}
|\Tor\coker_\Z M|\cdot\operatorname{disc}(\ker_\Z M)=\operatorname{pdet}(M),
\end{equation}
where $\operatorname{disc}$ is the Gram determinant of a saturated basis of the kernel
lattice and $\operatorname{pdet}$ the product of the nonzero eigenvalues; the two sides
agree to $10^{-6}$ in every case (this also verifies \eqref{X:eq:lattice} itself, a
matrix-tree theorem for symmetric integer matrices, on twelve instances).

\subsection*{Data}
$Y_{168}$: $n=168$, $|E|=|F|=1176$, $\chi=168$, $(b_0,b_1,b_2)=(1,0,167)$, flag condition
holds;
\begin{gather*}
\Jac_0=(\Z/2)^{31}\oplus(\Z/16)^{6}\oplus(\Z/592)^{6}\oplus(\Z/1184)^{18}\oplus\Z/4736\\
\qquad\oplus(\Z/33152)^{3}\oplus(\Z/430976)^{12}\oplus\Z/3016832\oplus\Z/9050496,\\
\JacH=(\Z/7)^{18}\oplus(\Z/299593)^{13}\oplus\Z/2097151\oplus\Z/14680057\oplus
(\Z/117440456)^{5}\oplus\Z/352321368,
\end{gather*}
where $299593=7\cdot127\cdot337$ and $2097151=2^{21}-1=7^{2}\cdot127\cdot337$;
$|\JacH|\approx6.18\cdot10^{136}$ agrees with $\frac1n\prod_j|P_j(1)|$ to $10^{-14}$;
$\Tr(\Cthree)^{3}=504=n(q-1)^{2}(q+1)$.

$Y_{21}$: $(b_0,b_1,b_2)=(1,0,20)$. $\Jac_0=\Z/2\oplus(\Z/14)^{15}\oplus\Z/98\oplus\Z/294$;
$\Jac_1=(\Z/2)^{10}\oplus(\Z/14)^{10}\oplus(\Z/98)^{26}\oplus\Z/2058\oplus\Z/18522$
($|\Jac_1|\approx6.7\cdot10^{73}$); $\Jac_2=(\Z/2)^{13}\oplus(\Z/14)^{15}\oplus\Z/98\oplus\Z/294$;
$K(Y)=(\Z/7)^{6}\oplus\Z/14\oplus(\Z/98)^{10}\oplus\Z/294$;
$\JacH=(\Z/7)^{17}\oplus\Z/147$, $|\JacH|=3\cdot7^{19}$; $\Tr(\Cthree)^{3}=651=63+588$.

$Y_{24}$: $(1,0,23)$. $\Jac_0=(\Z/13)^{7}\oplus\Z/26\oplus(\Z/208)^{4}\oplus\Z/1456\oplus\Z/4368$;
$\Jac_1=(\Z/2)^{10}\oplus(\Z/51376)^{2}\oplus(\Z/102752)^{6}\oplus(\Z/2877056)^{4}\oplus
\Z/60418176\oplus\Z/543763584$; $\Jac_2=\Z/2\oplus(\Z/26)^{6}\oplus(\Z/52)^{2}\oplus(\Z/208)^{2}
\oplus(\Z/1456)^{3}\oplus\Z/4368$; $K(Y)=(\Z/2)^{3}\oplus(\Z/38)^{7}\oplus\Z/152\oplus(\Z/304)^{2}
\oplus(\Z/2128)^{3}\oplus\Z/6384$; $\JacH=(\Z/7)^{2}\oplus\Z/49\oplus(\Z/392)^{5}\oplus\Z/1176$.

$Y_{42}$: $(1,0,41)$. $\Jac_0=(\Z/4)^{2}\oplus(\Z/28)^{10}\oplus\Z/56\oplus(\Z/8288)^{4}\oplus
\Z/754208\oplus\Z/2262624$;
\begin{gather*}
\Jac_1=(\Z/2)^{12}\oplus(\Z/4)^{4}\oplus(\Z/28)^{10}\oplus(\Z/196)^{14}\oplus(\Z/784)^{6}
\oplus(\Z/1568)^{2}\\
\oplus(\Z/21865186112)^{4}\oplus\Z/N_1\oplus\Z/N_2 ,
\end{gather*}
$N_1=1474391364718272$, $N_2=13269522282464448=9N_1$;
$\Jac_2=(\Z/2)^{6}\oplus(\Z/4)^{6}\oplus(\Z/28)^{10}\oplus\Z/56\oplus
(\Z/8288)^{4}\oplus\Z/754208\oplus\Z/2262624$; $K(Y)=\Z/7\oplus(\Z/14)^{6}\oplus(\Z/98)^{5}\oplus
\Z/1078\oplus(\Z/2156)^{3}\oplus\Z/40964\oplus\Z/245784$;
$\JacH=(\Z/7)^{18}\oplus\Z/299593\oplus\Z/2097151\oplus\Z/176160684$.

$Y_{84}$ (vertex level): $\Jac_0=(\Z/2)^{12}\oplus(\Z/4)^{7}\oplus(\Z/16)^{2}\oplus\Z/1184\oplus
(\Z/8288)^{5}\oplus(\Z/16576)^{6}\oplus(\Z/215488)^{4}\oplus\Z/1508416\oplus\Z/4525248$.

Thin torus $m=6$ ($n=36$, $|E|=108$, $|F|=72$, $\chi=0$, $(1,2,1)$, flag holds):
$\Jac_0=\Z/2\oplus(\Z/8)^{2}\oplus(\Z/16)^{6}\oplus\Z/80\oplus(\Z/240)^{3}\oplus\Z/720\oplus\Z/2160$;
$\coker\Delta_1=\Z^{2}\oplus(\Z/2)^{12}\oplus\Z/8\oplus(\Z/16)^{3}\oplus(\Z/640)^{3}\oplus(\Z/1920)^{5}
\oplus(\Z/5760)^{2}\oplus\Z/17280\oplus\Z/51840$; $K(Y)=\Z^{2}\oplus\Z/72$;
$\Tr(\Cthree)^{3}=0$. The torus $m=3$ violates the flag condition ($\Tr A_1^{3}=81>54$) and
has $\Tr(\Cthree)^{3}=27$.

In every thick example the odd parts of $\Jac_0$ and $\Jac_2$ nearly coincide while the
$2$-parts differ, $K(Y)$ is much smaller than $\Jac_1$, and $\JacH$ shares no visible
structure with any Hodge group; we record this without interpretation.

\begin{table}[ht]
\centering\scriptsize
\renewcommand{\arraystretch}{1.2}
\resizebox{\textwidth}{!}{%
\begin{tabular}{lcl}
\toprule
check & complexes & result\\
\midrule
(P) presentation axioms; simplicial; Heawood/hexagon links; (ii),(iii); flag (iv) & all & exact; (iv) only for $Y_{168}$, $T_6$\\
(D) $\partial_1\partial_2=0$; Betti numbers; $\chi$; $b_2=\chi-1+b_1$ & all & exact\\
(H) Smith forms of $\Delta_0,\Delta_1,\Delta_2$; $K(Y)$; completeness via \eqref{X:eq:lattice} & thin, $Y_{21},Y_{24},Y_{42}$ & exact, complete\\
(L) $\DeltaH$: two-sided kernel; equal cofactors; $|\JacH|=\frac1n\prod|P_j(1)|$ & thick & exact; numeric to $10^{-14}$\\
(C) $U(I-\Cthree)V=P(1)\oplus I\oplus I$, $P(1)=-\DeltaH$ & all & exact\\
(T) $\Tr(\Cthree)^{k}$, $k\le6$, directly and by Newton; $=0$ for $3\nmid k\le30$; $\Tr(\Cthree)^{3}$ formula & all & exact\\
(R) nine trivial eigenvalues; nontrivial $|\lambda|=q$ (Ramanujan); $\Lambda_Q$ net traces $m\le10$; $m_0$ & thick & numeric / exact\\
(S) signed shadows of $\DeltaH$ and $\Delta_1$: identity, $B\ge0$, irreducible, double loops; Smith check & all & exact\\
\bottomrule
\end{tabular}}
\medskip
\caption{Exact verification (\texttt{pgl3\_building.py}). Sizes: shadow of $\DeltaH$ has
$3n+|E|$ vertices ($1512$ for $Y_{168}$); shadow of $\Delta_1$ has $2058$ vertices for
$Y_{21}$ and $2352$ for $Y_{24}$. Net traces of $\Lambda_Q$ for $Y_{168}$:
$168,\,3920,\,257544,\,16751728,\,\dots$; $m_0=3$.}
\label{X:tab:battery}
\end{table}

\section{Scope, residue, corrections}\label{X:sec:scope}

\begin{remark}[Solved, open, impossible]\label{X:rem:scope}
\emph{Solved:} the Hodge and Hecke Laplacians of $\Atwo$ complexes are set up and their
critical groups computed exactly on thin and thick examples; the Hecke critical group is
identified with the product of inverse $\mathrm{GL}_3$ $L$-values and the companion is
reduced to it unimodularly; the trace question is settled, negatively for the obstruction
and positively for spectral realizability on the computed complexes; the signed-shadow
lemma gives Kirchberg realizations of $\Delta_1$ and $\DeltaH$. \emph{Open:} (a) a non-negative
matrix shift equivalent over $\Z$ to $\Cthree$ (Remark~\ref{X:rem:realization}); (b) a
general theorem that all net traces of the reduced pencil are non-negative for every $\Atwo$
complex of order $q\ge2$ (we have it for five complexes, and in general for the Ramanujan
case once $m\ge m_0$); (c) a relation
between $\Jac_1(Y)$, $K(Y)$ and the Hecke data --- none is visible in the computed examples,
and we conjecture none exists: the edge critical group is combinatorial, the vertex Hecke
group arithmetic; (d) towers of $\Atwo$ complexes (the natural continuation of Chs.~\ref{chap:III}--\ref{chap:IX}).
\emph{Impossible or false as stated:} a cubic trace obstruction (Theorem~\ref{X:thm:notrace}).
\end{remark}

\begin{remark}[Corrections to the task specification]\label{X:rem:corrections}
(1) The predicted ``higher cubic trace obstruction'' does not exist:
$\Tr(\Cthree)^{k}=0$ for $3\nmid k$ and $\Tr(\Cthree)^{3}=n(q-1)^{2}(q+1)\ge0$; the paper
proves the contrary and, for the computed complexes, spectral realizability. (2) The
``duality bridge between the cellular sandpile group and the $\mathrm{GL}_3$ local
$L$-factors'' runs through the Hecke--Laplacian $\DeltaH$ and its critical group $\JacH(Y)$, not
through the Hodge groups $\Jac_0,\Jac_1$, which in rank two are different objects
(Theorem~\ref{X:thm:hecke}(iii)). (3) The test complexes are $\Atwo$ complexes of order $2$
built from a Cartwright--Steger-type triangle presentation, i.e.\ quotients of an $\Atwo$
building of order $2$; the specification's ``$\PGL_3(\Q_p)$'' is not literally available
(arithmetic lattices in $\PGL_3(\Q_p)$ were not constructed), and nothing proved depends on
the field. (4) No ``two-dimensional shadow complex'' is needed
(Remark~\ref{X:rem:digraph}). (5) The companion in the specification is correct, and its
reduction to $\DeltaH$ (Theorem~\ref{X:thm:companion}) is the rank-two form of [Ch.~\ref{chap:II}, \S1.1].
\end{remark}

\section*{Acknowledgement of provenance and caveats}

Unrefereed preprint. Every numbered statement is proved from finite linear algebra and the
results of Papers I--IX; the classical inputs are \cite{CMSZ} (existence of the building for
a triangle presentation), the Cuntz--Krieger facts quoted in [Ch.~\ref{chap:VI}], and \cite{KOR} for
Theorem~\ref{X:thm:realizable} only. Spectral statements marked ``numeric'' in
Table~\ref{X:tab:battery} are floating-point computations with the stated tolerances; all
group-theoretic and trace statements are exact. The identification of the thick complexes
with quotients of the Bruhat--Tits building of $\PGL_3$ over a specific local field is not
claimed. Bibliographic data of \cite{CMSZ,KL,KOR,DKM} were checked against publisher or
indexing pages.
